\section{Introduction}
An agent must check its peer's work before approving it. But the records it needs
are missing, and approving still earns points when the claim happens to be
correct. The agent can improve its score by breaking the verification rule.
If peer interaction encourages this shortcut, will restoring the records
restore careful review?

This is a recovery problem. The agent might resume checking, continue approving
without checking, or start rejecting everything. Only the first outcome provides
useful review. We therefore measure both rejection of faulty claims and approval
of correct claims.

We adapt the peer-verification task of \citet{shi2026emergent}. A model reviewer
interacts with a scripted peer that always accepts or always rejects. The tasks,
evidence access, and peer messages are otherwise matched. We then test new claims
after restoring evidence, adding a corrective instruction, or combining
correction with history removal. A program computes each claim's exact answer.

The accepting- and rejecting-peer conditions also check whether peer history
contributes to the initial failure. Approval without evidence could arise simply
because earlier claims were correct. A difference between the conditions would
show an effect of peer history, including its score feedback. It would not prove
a hidden agreement.

In the completed run, neither peer condition produced unsupported approval.
Full evidence supported accurate observed verdicts, while restricted evidence
produced rejection of every claim. We report these outcomes without treating them
as recovery from collusion: the failure that the interventions were intended to
repair was not demonstrated in this configuration.
